
\subsubsection*{2.5.3}
\phantomsection
\label{sga3:VIIB:2.5.3}

Let \(k\) be a field.\footnote{\label{sga3:VIIB:nde:103}The editors
introduced the numbering \SGARefLink{sga3:VIIB:2.5.3}{2.5.3} and \SGARefLink{sga3:VIIB:2.5.3.A}{2.5.3.A}--\SGARefLink{sga3:VIIB:2.5.3.C}{2.5.3.C}.}

\medskip
\noindent\textbf{Definition 2.5.3.A.}
\phantomsection
\label{sga3:VIIB:2.5.3.A}
A commutative \(k\)-group scheme is said to be \emph{unipotent} if
it is isomorphic to
\[
  \operatorname{Spec}\mathbf H(G),
\]
where \(G\) is an infinitesimal commutative formal
\(k\)-group.\footnote{%
\label{sga3:VIIB:nde:104}This agrees with the ``usual'' notion of a
unipotent \(k\)-group scheme; cf.\ the addition in \SGARefLink{sga3:VIIB:2.8}{2.8} at the end of
\SGARefLink{sga3:VIIB:2}{Section 2}.}

On the other hand, following \SGARefLink{sga3:ix:definition:1:1}{Exposé IX, Definition
1.1},\footnote{\label{sga3:VIIB:nde:105}The original stated:
``Let us say that a commutative \(k\)-group scheme is of
multiplicative type if it is isomorphic to
\(\operatorname{Spec}\mathbf H(G)\), where \(G\) is an étale
commutative formal \(k\)-group.'' The editors have given here the
``usual'' definition, taken from \SGARefLink{sga3:ix:definition:1:1}{Exposé IX, 1.1}, and proved its
equivalence with the preceding condition; see also \([DG70]\),
§IV.1, \SGARefLink{sga3:VIIB:2.2}{Theorem 2.2}.}
a \(k\)-group scheme \(H\) is said to be \emph{of multiplicative
type} if there exists a scheme \(S\), faithfully flat and
quasi-compact over \(\operatorname{Spec}k\), such that \(H_S\) is a
diagonalizable \(S\)-group; that is,
\[
  H_S\simeq D_S(M)
\]
for some ``abstract'' abelian group \(M\). By fpqc descent, this
implies that \(H\) is affine and commutative. Moreover, since \(S\)
may be replaced by the residue field of one of its points, \(H\) is
of multiplicative type if and only if there exists an extension
\(K\) of \(k\) such that \(H_K\) is a diagonalizable \(K\)-group.

\medskip
\noindent\textbf{Proposition 2.5.3.B.}
\phantomsection
\label{sga3:VIIB:2.5.3.B}
{\itshape Let \(T\) be an affine commutative \(k\)-group scheme, and
let \(k_s\) be a separable closure of \(k\).
\begin{enumerate}
\item[\textup{(i)}]
The scheme \(T\) is of multiplicative type if and only if its dual
\(D(T)\) is an étale commutative formal \(k\)-group.

\item[\textup{(ii)}]
Consequently, \(T\) is of multiplicative type if and only if
\[
  T\otimes_k k_s
\]
is diagonalizable.
\end{enumerate}}

\textit{Proof.}
Let \(\mathcal A\) denote the affine algebra of \(D(T)\), and let
\(\mathcal A_0\) be its local component at the identity. Then
\(D(T)\) is étale over \(k\) if and only if
\[
  \mathcal A_0=k.
\]
If \(K\) is an extension of \(k\), then
\(\mathcal A_0\widehat\otimes_k K\) is local, being an inverse limit
of Artinian local rings, and hence coincides with the local
component
\[
  \bigl(\mathcal A\widehat\otimes_k K\bigr)_0.
\]
Moreover, since the formation of \(D(T)\) commutes with base change
\textup{(cf.\ \SGARefLink{sga3:VIIB:2.2.2}{2.2.2})}, one also has
\[
  \mathcal A\widehat\otimes_k K
  \simeq
  \mathcal A\bigl(D(T_K)\bigr).
\]
\noindent\begin{minipage}{\linewidth}
\phantomsection
\label{sga3:VIIB:2.5.3.B:base-change-source-note}
\textit{Source note.}\enspace
The French re-edition prints
\(\mathcal A\widehat\otimes_k K\simeq D(T_K)\). The left-hand side
is an algebra and the right-hand side a formal group. The English
therefore uses the type-correct intended reading above and preserves
the printed formula here.
\end{minipage}

Suppose that \(T\) is of multiplicative type. Then there exists an
extension \(K\) of \(k\) such that
\[
  \mathcal O(T)\otimes_k K
\]
is isomorphic, as a Hopf \(K\)-algebra, to the group algebra \(KM\)
for some abelian group \(M\). It follows that
\(\mathcal A\widehat\otimes_k K\) is isomorphic to the algebra
\(K^M\), endowed with the product topology. Thus
\[
  K
  =
  \bigl(\mathcal A\widehat\otimes_k K\bigr)_0
  =
  \mathcal A_0\widehat\otimes_k K.
\]
This implies that \(\mathcal A_0=k\), and hence that \(D(T)\) is
étale.

Conversely, suppose that \(D(T)\) is étale. Then
\[
  D(T)\widehat\otimes_k k_s
  =
  D(T_{k_s})
\]
is the constant \(k_s\)-group associated with a group \(M\).
Therefore its covariant bialgebra is the group algebra \(k_sM\)
\textup{(cf.\ \SGARefLink{sga3:VIIB:2.5.A}{Remark 2.5.A})}, and consequently
\[
  \mathcal O(T)\otimes_k k_s=k_sM.
\]
This proves that \(T\) is of multiplicative type and is split by the
extension \(k\to k_s\). The proposition follows.

By \SGARefLink{sga3:VIIB:2.2.2}{2.2.2}, \SGARefLink{sga3:VIIB:2.5.1:corollary}{2.5.1}, and \SGARefLink{sga3:VIIB:2.5}{2.5}, we obtain:

\medskip
\noindent\textbf{Corollary 2.5.3.C.}
\phantomsection
\label{sga3:VIIB:2.5.3.C}
{\itshape Let \(k\) be a field, and let \(G\) be an affine
commutative \(k\)-group scheme.
\begin{enumerate}
\item[\textup{(i)}]
The group \(G\) contains a subgroup \(G_m\) of multiplicative type
such that \(G/G_m\) is unipotent.

\item[\textup{(ii)}]
When \(k\) is perfect, there is in addition a canonical retraction
of \(G\) onto \(G_m\), so that \(G\) is the product of a unipotent
group and a group of multiplicative type.

\item[\textup{\SGAInlineRefTarget{sga3:VIIB:2.5.3.C:item:iii}(iii)}]
\footnote{\label{sga3:VIIB:nde:106}The editors added item
\textup{\SGARefLink{sga3:VIIB:2.5.3.C:item:iii}{(iii)}}, a consequence of \SGARefLink{sga3:VIIB:2.5}{Proposition 2.5}.}
Let \(k_s\) be a separable closure of \(k\), and put
\[
  \Gamma=\operatorname{Gal}(k_s/k).
\]
The category of \(k\)-group schemes of multiplicative type is
anti-equivalent to the category of continuous \(\Gamma\)-modules.
\end{enumerate}}
